\cotefile{156-4}
\scanpage{66}\mspage{65}
Ceci résulte du \deleted{\uline{C}4 et} cor. 1 et du

\uline{Lemme} Pour que $F$ soit de t. f.\inserted{(resp. \unsure{fini|finie})} Il f et s que a) les $F_i$ le soient et b) Tous les $F_i$ sauf un \unsure{nb} fini sont $\emptyset$.

Notations $\pi_0(F)$ (NB Si $F\leq G$ on a $\pi_0(F)\longrightarrow\pi_0(G)$\deleted{\unsure{)|?}}, et transitivité) $\pi_0(M)$, $\pi_0^{\deleted{\unsure{\mathrm{cat}|\mathrm{act}}}!}(\mathcal F)$. (\unsure{Attention}, le $\pi_0^{\unsure{d|!}}(\mathcal F)$ n'est pas le $\pi_0(F,\leq)$, mais le $\pi_0$ ``en tenant compte \unsure{du|des} \unsure{Sup}''. Si $\unsure{\mathcal C|\mathscr C|C}=(\mathcal F,\leq,\deleted{\unsure{\leq|\prec}}\unsure{\preceq|\leqslant})$\margin{\unsure{Contrée|Contrie}}\margin{qui ne sert \deleted{\unsure{pourtant}} \unsure{à} rien ici}, on pose aussi $\pi_0(\mathcal C)$.

Axiome renforcé

\uuline{C4 bis} Pour tout $(F_i)_i\in\prod\mathcal F_i$, $\deleted{\unsure{\exists F|\mathcal F}}$ $\exists F\in\mathcal F$ telle que $F\mid M_i=\deleted{\unsure{\mathcal F}}F_i$ pour $i\unsure{\vdash|\in}I$.

\uline{\unsure{Contrée|Contrie} \unsure{résiduelle|réssiduelle}} \uline{Support} \uline{et} \uline{cosupport}

Soit $\deleted{\unsure{(F_i)\in\mathcal F|(H_i)\in\mathcal F}}$ $\Phi\subset\mathcal F$ un ens. de figures \inserted{$H_i$}. \deleted{On posera} Soit $F$ une figure. Pour \deleted{\unsure{Cosupp.} $\Phi$} que $F$ soit \unsure{disjointe|disjoint} de tous les $H_i$, il f et s. que ce soit le cas pour toutes les strates de $F$ (à cause de C3). \deleted{\unsure{Posons|Poser}} Posons
